\chapter{ Discussion }

% The enabling of third-party plugins on embedded audio systems is highly desirable, but blah blah blah.

Performance benchmarking in Chapter~\ref{sec:bench} demonstrates that WebAssembly (WASM) can support real-time audio DSP on embedded hardware under certain conditions, but further research is needed to determine the bounds of its performance.

Across the four tested algorithms, interpreter-based execution (\lstinline{wasmi}) did not achieve usable throughput and can be discarded as a runtime approach in future work. Ahead-of-time (AOT) compilation using WAMR and translation via \lstinline{wasm2c} both yielded viable performance for algorithms whose memory footprint is known at compile time, but demonstrated highly variable success with algorithms that do not.

Two patterns are evident. First, strategies that produce native machine code at flash time are markedly more suitable for real-time audio than runtime interpretation of bytecode. Second, runtime overhead is highly workload-dependent. Algorithms that rely heavily on floating-point arithmetic but maintain predictable control flow and limited dynamic allocation are more amenable to WASM deployment.

It is important to note that these measurements were obtained without aggressive runtime customization. Further opportunities for performance gains exist, including better leveraging of host-imported math intrinsics, improved compiler pipelines, and tighter ABI and runtime specialization. The results presented in this work therefore represent a lower bound on achievable performance.

% Even with the limitations displayed in this limited benchmark, the affordances of a universal distributable for audio DSP may in many cases outweigh the performance drawbacks.

A reasonable question remains: if real-time performance of WASM-based DSP modules relies on AOT translation, can they truly be considered a universal distributable? The author argues yes. I encourage conception of the AOT compiler as part of the runtime or flash tool for each target, a responsibility laid on hardware manufacturers looking to increase adoption of their product rather than a burden on audio plugin developers. In some cases, AOT compilation for embedded systems can even occur on the target device, as seen in \cite{10756513}.

% While dependence on an AOT compiler for each target

% may be considered akin to an installer. On-device AOT shrinks this distinction further.

% on-device AOT compilation as discussed in \cite{10756513}.
